\Subhead{1.5.2}{Вынужденные колебания}
Гармоническая сила возбуждает нормальные моды через пространственную проекцию. Для модели $M\ddot{\bm q}+C\dot{\bm q}+K\bm q=\bm F$ получаем следующий отклик. \paragraph{Вынужденный отклик и потери.}
С учётом вязкой матрицы потерь $C$ гармоническая амплитуда равна
\begin{equation}
 \bm q_\Omega=(K-\Omega^2M-\ii\Omega C)^{-1}\bm F_\Omega.
\end{equation}
Если $\bm a_m^{\mathsf T}C\bm a_n=2\gamma_n\delta_{mn}$, то
\begin{equation}
 Q_{n,\Omega}=\frac{\bm a_n^{\mathsf T}\bm F_\Omega}
 {\omega_n^2-\Omega^2-2\ii\gamma_n\Omega}.
\end{equation}
Проекция источника определяет, будет ли мода возбуждена: при нулевой проекции резонанс не наблюдается даже на собственной частоте. Для произвольных непропорциональных потерь $C$ модальные уравнения остаются связанными; диагонализация только $M,K$ не позволяет приписать каждой моде независимое трение. Между резонансами возможны антирезонансы: отдельная измеряемая координата обращается в нуль вследствие компенсации модальных вкладов, хотя система продолжает двигаться.

Для двух одинаковых осцилляторов синфазную моду возбуждает сумма сил $F_1+F_2$, противофазную --- разность $F_1-F_2$. Поэтому два одинаковых синфазных воздействия не возбуждают противофазную моду даже на её собственной частоте. Так проявляется пространственное правило отбора.
